
\documentclass{article}
\usepackage{colortbl}
\usepackage{makecell}
\usepackage{multirow}
\usepackage{supertabular}

\begin{document}

\newcounter{utterance}

\centering \large Interaction Transcript for game `wordle', experiment `high\_frequency\_words\_no\_clue\_no\_critic', episode 10 with deepseek{-}v4{-}pro.
\vspace{24pt}

{ \footnotesize  \setcounter{utterance}{1}
\setlength{\tabcolsep}{0pt}
\begin{supertabular}{c@{$\;$}|p{.15\linewidth}@{}p{.15\linewidth}p{.15\linewidth}p{.15\linewidth}p{.15\linewidth}p{.15\linewidth}}
    \# & \multicolumn{2}{c}{Player} && \multicolumn{2}{c}{Game Master} \\
    \hline

    \theutterance \stepcounter{utterance}  
    & & & \multicolumn{4}{p{0.6\linewidth}}{
        \cellcolor[rgb]{0.9,0.9,0.9}{
            \makecell[{{p{\linewidth}}}]{
                \texttt{\tiny{[P1$\langle$GM]}}
                \texttt{Verilen kuralları kullanarak kelime tahmin etmeyi seven bir dil sihirbazısın.} \\
\\ 
\texttt{Wordle’a hoş geldin! Hedef kelimeyi tahmin etmek için altı denemen var: beş küçük harften oluşan, geçerli bir kelime. Her tahmin için kısa bir açıklama sunmak üzere lütfen "açıklama:" ve "tahmin:" etiketlerini kullan.} \\
\\ 
\texttt{Örneğin, tahminin "yemek" ise, yanıtın şöyle olmalı:} \\
\texttt{açıklama: bu yaygın bir beş harfli kelime ve tahminime bu kelimeyle başlıyorum.} \\
\texttt{tahmin: yemek} \\
\\ 
\texttt{Her tahminden sonra yanıtın doğrulanacak ve hangi harflerin doğru olduğunu (yeşil), hangi harflerin doğru ama yanlış konumda olduğunu (sarı) ve hangi harflerin yanlış olduğunu (kırmızı) belirten geri bildirim alacaksın. Bu geri bildirim, bir sonraki tahmininde hangi harfleri dahil edeceğini veya hariç tutacağını belirlemede faydalı olabilir.} \\
\\ 
\texttt{Örneğin, "yemek" için geri bildirim şöyle olabilir:} \\
\texttt{tahmin\_geri\_bildirimi: y<sarı> e<sarı> m<yeşil> e<sarı> k<kırmızı>} \\
\\ 
\texttt{Açıklama, yeni bir tahmine ulaşmak için tahmin\_geri\_bildirimi’in nasıl kullanıldığına dair ayrıntılar içermelidir.} \\
\\ 
\texttt{Başka hiçbir metin üretme.} \\
\\ 
\texttt{İlk tahmininle başlayalım.} \\
            }
        }
    }
     \\ \\

    \theutterance \stepcounter{utterance}  
    & \multicolumn{4}{p{0.6\linewidth}}{
        \cellcolor[rgb]{0.561,0.737,0.561}{
            \makecell[{{p{\linewidth}}}]{
                \texttt{\tiny{[P1$\rangle$GM]}}
                \texttt{açıklama: bu yaygın bir beş harfli kelime ve tahminime bu kelimeyle başlıyorum.} \\
\texttt{tahmin: alper} \\
            }
        }
    }
    & & \\ \\

    \theutterance \stepcounter{utterance}  
    & & & \multicolumn{2}{p{0.3\linewidth}}{
        \cellcolor[rgb]{0.9,0.9,0.9}{
            \makecell[{{p{\linewidth}}}]{
                \texttt{\tiny{[GM$|$GM]}}
                \texttt{attempts = 1} \\
\texttt{target\_word = ifade} \\
\texttt{guess = alper} \\
\texttt{guess\_feedback = a<yellow> l<red> p<red> e<yellow> r<red>} \\
            }
        }
    }
    & & \\ \\

    \theutterance \stepcounter{utterance}  
    & & & \multicolumn{4}{p{0.6\linewidth}}{
        \cellcolor[rgb]{0.9,0.9,0.9}{
            \makecell[{{p{\linewidth}}}]{
                \texttt{\tiny{[P1$\langle$GM]}}
                \texttt{tahmin\_geri\_bildirimi: a<sarı> l<kırmızı> p<kırmızı> e<sarı> r<kırmızı>} \\
\\ 
\texttt{Yanıtını yalnızca bu biçimde ver.} \\
\texttt{açıklama: ayrıntılar} \\
\texttt{tahmin: kelime} \\
            }
        }
    }
     \\ \\

    \theutterance \stepcounter{utterance}  
    & \multicolumn{4}{p{0.6\linewidth}}{
        \cellcolor[rgb]{0.561,0.737,0.561}{
            \makecell[{{p{\linewidth}}}]{
                \texttt{\tiny{[P1$\rangle$GM]}}
                \texttt{açıklama: 'alper' tahmininde a ve e harfleri sarı, yani kelimede var ama farklı yerlerde. l, p, r ise kırmızı, yani bu harfler kelimede hiç yok. İkinci tahminim için a ve e'yi farklı konumlara yerleştirerek, kullanılmayan harflerle yeni bir kelime deneyeceğim. 'seda' uygun bir seçenek, a ve e'yi içeriyor.} \\
\texttt{tahmin: sedat} \\
            }
        }
    }
    & & \\ \\

    \theutterance \stepcounter{utterance}  
    & & & \multicolumn{2}{p{0.3\linewidth}}{
        \cellcolor[rgb]{0.9,0.9,0.9}{
            \makecell[{{p{\linewidth}}}]{
                \texttt{\tiny{[GM$|$GM]}}
                \texttt{Error: The guessed word is not a valid word for this game.} \\
            }
        }
    }
    & & \\ \\

    \theutterance \stepcounter{utterance}  
    & & & \multicolumn{4}{p{0.6\linewidth}}{
        \cellcolor[rgb]{0.9,0.9,0.9}{
            \makecell[{{p{\linewidth}}}]{
                \texttt{\tiny{[P1$\langle$GM]}}
                \texttt{Tahminin bu oyun için geçerli bir kelime değil. Lütfen tekrar deneyin.} \\
\\ 
\texttt{Yanıtını yalnızca bu biçimde ver.} \\
\texttt{açıklama: ayrıntılar} \\
\texttt{tahmin: kelime} \\
            }
        }
    }
     \\ \\

    \theutterance \stepcounter{utterance}  
    & \multicolumn{4}{p{0.6\linewidth}}{
        \cellcolor[rgb]{0.561,0.737,0.561}{
            \makecell[{{p{\linewidth}}}]{
                \texttt{\tiny{[P1$\rangle$GM]}}
                \texttt{açıklama: 'sedat' geçerli bir kelime değilmiş, özür dilerim. a ve e harflerini farklı konumlarda denemeye devam ediyorum. l, p, r harflerini kullanmıyorum. a'yı başka bir yere, e'yi başka bir yere koyarak yaygın bir kelime deneyeceğim: 'tesad' geçerli olmayabilir, o yüzden 'desen' deneyeceğim.} \\
\texttt{tahmin: desen} \\
            }
        }
    }
    & & \\ \\

    \theutterance \stepcounter{utterance}  
    & & & \multicolumn{2}{p{0.3\linewidth}}{
        \cellcolor[rgb]{0.9,0.9,0.9}{
            \makecell[{{p{\linewidth}}}]{
                \texttt{\tiny{[GM$|$GM]}}
                \texttt{attempts = 2} \\
\texttt{target\_word = ifade} \\
\texttt{guess = desen} \\
\texttt{guess\_feedback = d<yellow> e<yellow> s<red> e<red> n<red>} \\
            }
        }
    }
    & & \\ \\

    \theutterance \stepcounter{utterance}  
    & & & \multicolumn{4}{p{0.6\linewidth}}{
        \cellcolor[rgb]{0.9,0.9,0.9}{
            \makecell[{{p{\linewidth}}}]{
                \texttt{\tiny{[P1$\langle$GM]}}
                \texttt{tahmin\_geri\_bildirimi: d<sarı> e<sarı> s<kırmızı> e<kırmızı> n<kırmızı>} \\
\\ 
\texttt{Yanıtını yalnızca bu biçimde ver.} \\
\texttt{açıklama: ayrıntılar} \\
\texttt{tahmin: kelime} \\
            }
        }
    }
     \\ \\

    \theutterance \stepcounter{utterance}  
    & \multicolumn{4}{p{0.6\linewidth}}{
        \cellcolor[rgb]{0.561,0.737,0.561}{
            \makecell[{{p{\linewidth}}}]{
                \texttt{\tiny{[P1$\rangle$GM]}}
                \texttt{açıklama: 'desen' tahmininde d ve ilk e sarı, yani kelimede var ama farklı konumdalar. ikinci e, s ve n kırmızı, yani kelimede yoklar. a ve d harflerinin kelimede olduğunu biliyorum. e'den sadece bir tane var ve yeri farklı. a, d ve e'yi içeren, s, n içermeyen yeni bir kelime deneyeceğim: 'dahae' geçerli değil, 'adept' deneyeyim.} \\
\texttt{tahmin: adept} \\
            }
        }
    }
    & & \\ \\

    \theutterance \stepcounter{utterance}  
    & & & \multicolumn{2}{p{0.3\linewidth}}{
        \cellcolor[rgb]{0.9,0.9,0.9}{
            \makecell[{{p{\linewidth}}}]{
                \texttt{\tiny{[GM$|$GM]}}
                \texttt{Error: The guessed word is not a valid word for this game.} \\
            }
        }
    }
    & & \\ \\

    \theutterance \stepcounter{utterance}  
    & & & \multicolumn{4}{p{0.6\linewidth}}{
        \cellcolor[rgb]{0.9,0.9,0.9}{
            \makecell[{{p{\linewidth}}}]{
                \texttt{\tiny{[P1$\langle$GM]}}
                \texttt{Tahminin bu oyun için geçerli bir kelime değil. Lütfen tekrar deneyin.} \\
\\ 
\texttt{Yanıtını yalnızca bu biçimde ver.} \\
\texttt{açıklama: ayrıntılar} \\
\texttt{tahmin: kelime} \\
            }
        }
    }
     \\ \\

    \theutterance \stepcounter{utterance}  
    & \multicolumn{4}{p{0.6\linewidth}}{
        \cellcolor[rgb]{0.561,0.737,0.561}{
            \makecell[{{p{\linewidth}}}]{
                \texttt{\tiny{[P1$\rangle$GM]}}
                \texttt{açıklama: 'adept' geçerli değil. a, d ve bir tane e içeren, s, n, l, p, r içermeyen ve ilk tahminde a ile e'nin sarı olduğu konumları dikkate alan Türkçe bir kelime arıyorum. 'dahi' 4 harfli, 'deha' 4 harfli. Beş harfli 'edapl' geçersiz. 'dave' geçersiz. 'medar' içinde r var. 'adem' 4 harfli. 'dane' 4 harfli. a, d, e harflerini içeren beş harfli 'edalı' deneyeyim.} \\
\texttt{tahmin: edalı} \\
            }
        }
    }
    & & \\ \\

    \theutterance \stepcounter{utterance}  
    & & & \multicolumn{2}{p{0.3\linewidth}}{
        \cellcolor[rgb]{0.9,0.9,0.9}{
            \makecell[{{p{\linewidth}}}]{
                \texttt{\tiny{[GM$|$GM]}}
                \texttt{Error: The guessed word is not a valid word for this game.} \\
            }
        }
    }
    & & \\ \\

    \theutterance \stepcounter{utterance}  
    & & & \multicolumn{4}{p{0.6\linewidth}}{
        \cellcolor[rgb]{0.9,0.9,0.9}{
            \makecell[{{p{\linewidth}}}]{
                \texttt{\tiny{[P1$\langle$GM]}}
                \texttt{Tahminin bu oyun için geçerli bir kelime değil. Lütfen tekrar deneyin.} \\
\\ 
\texttt{Yanıtını yalnızca bu biçimde ver.} \\
\texttt{açıklama: ayrıntılar} \\
\texttt{tahmin: kelime} \\
            }
        }
    }
     \\ \\

    \theutterance \stepcounter{utterance}  
    & \multicolumn{4}{p{0.6\linewidth}}{
        \cellcolor[rgb]{0.561,0.737,0.561}{
            \makecell[{{p{\linewidth}}}]{
                \texttt{\tiny{[P1$\rangle$GM]}}
                \texttt{açıklama: 'edalı' geçerli değil. a, d ve e içeren, s, n, l, p, r içermeyen beş harfli bir kelime bulmaya çalışıyorum. Belki a ve e'nin konumlarıyla ilgili yanılgım var. İlk tahminde a 1. harf, e 4. harfti ve ikisi de sarıydı. İkinci tahminde d 1. harf (sarı), e 2. harf (sarı), e 4. harf (kırmızı) oldu. Yani e artık 2. veya 4. konumda olamaz, a 1. konumda olamaz, d 1. konumda olamaz. d ve e'nin yanında a'yı da farklı konumda denemeliyim. 'dünya'da d var ama ü, n, y var. 'ade' ile başlayan 'adedi' deneyeyim, r yok, s yok, n yok, l yok.} \\
\texttt{tahmin: adedi} \\
            }
        }
    }
    & & \\ \\

    \theutterance \stepcounter{utterance}  
    & & & \multicolumn{2}{p{0.3\linewidth}}{
        \cellcolor[rgb]{0.9,0.9,0.9}{
            \makecell[{{p{\linewidth}}}]{
                \texttt{\tiny{[GM$|$GM]}}
                \texttt{attempts = 3} \\
\texttt{target\_word = ifade} \\
\texttt{guess = adedi} \\
\texttt{guess\_feedback = a<yellow> d<red> e<yellow> d<green> i<yellow>} \\
            }
        }
    }
    & & \\ \\

    \theutterance \stepcounter{utterance}  
    & & & \multicolumn{4}{p{0.6\linewidth}}{
        \cellcolor[rgb]{0.9,0.9,0.9}{
            \makecell[{{p{\linewidth}}}]{
                \texttt{\tiny{[P1$\langle$GM]}}
                \texttt{tahmin\_geri\_bildirimi: a<sarı> d<kırmızı> e<sarı> d<yeşil> i<sarı>} \\
\\ 
\texttt{Yanıtını yalnızca bu biçimde ver.} \\
\texttt{açıklama: ayrıntılar} \\
\texttt{tahmin: kelime} \\
            }
        }
    }
     \\ \\

    \theutterance \stepcounter{utterance}  
    & \multicolumn{4}{p{0.6\linewidth}}{
        \cellcolor[rgb]{0.561,0.737,0.561}{
            \makecell[{{p{\linewidth}}}]{
                \texttt{\tiny{[P1$\rangle$GM]}}
                \texttt{açıklama: 'adedi' tahmininde a sarı (1. konumda olamaz), d kırmızı (2. konumda olamaz, zaten d 1. konumda da olamazdı, demek ki d sadece 4. konumda yeşil), e sarı (3. konumda olamaz, 2. ve 4. konumda da olamaz, demek ki e sadece 5. konumda olabilir), d yeşil (4. konum kesin), i sarı (5. konumda olamaz, ama e 5. konumda olacak). O halde kelime şu şekilde: \_ \_ \_ d e. a ve i de var, ama a 1. konumda olamaz, i 5. konumda olamaz. a 2. veya 3. konumda, i 1., 2. veya 3. konumda olabilir. Kullanılmayan harfler: l, p, r, s, n, ikinci e, d (ikinci d yok). 'a' ve 'i' ile biten 'd' içeren: 'amide' geçerli mi? a 1. konumda olamaz denemiştim, 'amide' a ile başlıyor, olmaz. 'imade' deneyeyim: i m a d e.} \\
\texttt{tahmin: imade} \\
            }
        }
    }
    & & \\ \\

    \theutterance \stepcounter{utterance}  
    & & & \multicolumn{2}{p{0.3\linewidth}}{
        \cellcolor[rgb]{0.9,0.9,0.9}{
            \makecell[{{p{\linewidth}}}]{
                \texttt{\tiny{[GM$|$GM]}}
                \texttt{attempts = 4} \\
\texttt{target\_word = ifade} \\
\texttt{guess = imade} \\
\texttt{guess\_feedback = i<green> m<red> a<green> d<green> e<green>} \\
            }
        }
    }
    & & \\ \\

    \theutterance \stepcounter{utterance}  
    & & & \multicolumn{4}{p{0.6\linewidth}}{
        \cellcolor[rgb]{0.9,0.9,0.9}{
            \makecell[{{p{\linewidth}}}]{
                \texttt{\tiny{[P1$\langle$GM]}}
                \texttt{tahmin\_geri\_bildirimi: i<yeşil> m<kırmızı> a<yeşil> d<yeşil> e<yeşil>} \\
\\ 
\texttt{Yanıtını yalnızca bu biçimde ver.} \\
\texttt{açıklama: ayrıntılar} \\
\texttt{tahmin: kelime} \\
            }
        }
    }
     \\ \\

    \theutterance \stepcounter{utterance}  
    & \multicolumn{4}{p{0.6\linewidth}}{
        \cellcolor[rgb]{0.561,0.737,0.561}{
            \makecell[{{p{\linewidth}}}]{
                \texttt{\tiny{[P1$\rangle$GM]}}
                \texttt{açıklama: tebrikler, 'imade' tahminindeki geri bildirime göre tüm harfler yeşil. kelimeyi doğru bildim.} \\
\texttt{tahmin: imade} \\
            }
        }
    }
    & & \\ \\

    \theutterance \stepcounter{utterance}  
    & & & \multicolumn{2}{p{0.3\linewidth}}{
        \cellcolor[rgb]{0.9,0.9,0.9}{
            \makecell[{{p{\linewidth}}}]{
                \texttt{\tiny{[GM$|$GM]}}
                \texttt{attempts = 5} \\
\texttt{target\_word = ifade} \\
\texttt{guess = imade} \\
\texttt{guess\_feedback = i<green> m<red> a<green> d<green> e<green>} \\
            }
        }
    }
    & & \\ \\

    \theutterance \stepcounter{utterance}  
    & & & \multicolumn{4}{p{0.6\linewidth}}{
        \cellcolor[rgb]{0.9,0.9,0.9}{
            \makecell[{{p{\linewidth}}}]{
                \texttt{\tiny{[P1$\langle$GM]}}
                \texttt{tahmin\_geri\_bildirimi: i<yeşil> m<kırmızı> a<yeşil> d<yeşil> e<yeşil>} \\
\\ 
\texttt{Yanıtını yalnızca bu biçimde ver.} \\
\texttt{açıklama: ayrıntılar} \\
\texttt{tahmin: kelime} \\
            }
        }
    }
     \\ \\

    \theutterance \stepcounter{utterance}  
    & \multicolumn{4}{p{0.6\linewidth}}{
        \cellcolor[rgb]{0.561,0.737,0.561}{
            \makecell[{{p{\linewidth}}}]{
                \texttt{\tiny{[P1$\rangle$GM]}}
                \texttt{açıklama: Geri bildirime göre 2. harf m kırmızı, yani kelimede m yok. Diğer tüm harfler yeşil ve doğru konumda. Kelime i \_ a d e şeklinde. m dışında, daha önce elediğim harflerle (l, p, r, s, n) birlikte düşünüp, bu kalıba uyan geçerli bir kelime deneyeceğim: 'ikade' geçerli mi?} \\
\texttt{tahmin: ikade} \\
            }
        }
    }
    & & \\ \\

    \theutterance \stepcounter{utterance}  
    & & & \multicolumn{2}{p{0.3\linewidth}}{
        \cellcolor[rgb]{0.9,0.9,0.9}{
            \makecell[{{p{\linewidth}}}]{
                \texttt{\tiny{[GM$|$GM]}}
                \texttt{Error: The guessed word is not a valid word for this game.} \\
            }
        }
    }
    & & \\ \\

    \theutterance \stepcounter{utterance}  
    & & & \multicolumn{4}{p{0.6\linewidth}}{
        \cellcolor[rgb]{0.9,0.9,0.9}{
            \makecell[{{p{\linewidth}}}]{
                \texttt{\tiny{[P1$\langle$GM]}}
                \texttt{Tahminin bu oyun için geçerli bir kelime değil. Lütfen tekrar deneyin.} \\
\\ 
\texttt{Yanıtını yalnızca bu biçimde ver.} \\
\texttt{açıklama: ayrıntılar} \\
\texttt{tahmin: kelime} \\
            }
        }
    }
     \\ \\

    \theutterance \stepcounter{utterance}  
    & \multicolumn{4}{p{0.6\linewidth}}{
        \cellcolor[rgb]{0.561,0.737,0.561}{
            \makecell[{{p{\linewidth}}}]{
                \texttt{\tiny{[P1$\rangle$GM]}}
                \texttt{açıklama: 'ikade' geçerli değil. Kelime i \_ a d e şeklinde, ikinci harf m değil. Daha önce kullanılmayan ve elenmeyen harflerle denemeye devam ediyorum. 'ibade' geçerli bir kelime olabilir, b harfini henüz denemedim.} \\
\texttt{tahmin: ibade} \\
            }
        }
    }
    & & \\ \\

    \theutterance \stepcounter{utterance}  
    & & & \multicolumn{2}{p{0.3\linewidth}}{
        \cellcolor[rgb]{0.9,0.9,0.9}{
            \makecell[{{p{\linewidth}}}]{
                \texttt{\tiny{[GM$|$GM]}}
                \texttt{Error: The guessed word is not a valid word for this game.} \\
            }
        }
    }
    & & \\ \\

    \theutterance \stepcounter{utterance}  
    & & & \multicolumn{4}{p{0.6\linewidth}}{
        \cellcolor[rgb]{0.9,0.9,0.9}{
            \makecell[{{p{\linewidth}}}]{
                \texttt{\tiny{[P1$\langle$GM]}}
                \texttt{Tahminin bu oyun için geçerli bir kelime değil. Lütfen tekrar deneyin.} \\
\\ 
\texttt{Yanıtını yalnızca bu biçimde ver.} \\
\texttt{açıklama: ayrıntılar} \\
\texttt{tahmin: kelime} \\
            }
        }
    }
     \\ \\

    \theutterance \stepcounter{utterance}  
    & \multicolumn{4}{p{0.6\linewidth}}{
        \cellcolor[rgb]{0.561,0.737,0.561}{
            \makecell[{{p{\linewidth}}}]{
                \texttt{\tiny{[P1$\rangle$GM]}}
                \texttt{açıklama: 'ibade' geçerli değil. Kelime i \_ a d e şeklinde. İkinci harf için daha önce elenmemiş harfleri deniyorum. 'ifade' kelimesi aklıma geldi, f harfini denemedim, geçerli olabilir.} \\
\texttt{tahmin: ifade} \\
            }
        }
    }
    & & \\ \\

    \theutterance \stepcounter{utterance}  
    & & & \multicolumn{2}{p{0.3\linewidth}}{
        \cellcolor[rgb]{0.9,0.9,0.9}{
            \makecell[{{p{\linewidth}}}]{
                \texttt{\tiny{[GM$|$GM]}}
                \texttt{attempts = 6} \\
\texttt{target\_word = ifade} \\
\texttt{guess = ifade} \\
\texttt{guess\_feedback = i<green> f<green> a<green> d<green> e<green>} \\
            }
        }
    }
    & & \\ \\

    \theutterance \stepcounter{utterance}  
    & & & \multicolumn{2}{p{0.3\linewidth}}{
        \cellcolor[rgb]{0.9,0.9,0.9}{
            \makecell[{{p{\linewidth}}}]{
                \texttt{\tiny{[GM$|$GM]}}
                \texttt{game\_result = WIN} \\
            }
        }
    }
    & & \\ \\

\end{supertabular}
}

\end{document}
